\pdfinfoomitdate=1
\pdftrailerid{}
\pdfsuppressptexinfo=15
\documentclass[11pt,letterpaper]{article}
\usepackage[margin=0.88in]{geometry}
\usepackage[T1]{fontenc}
\usepackage{lmodern,microtype,amsmath,amssymb,amsthm,mathtools,booktabs}
\usepackage[hidelinks]{hyperref}
\usepackage{enumitem}
\setlist{nosep}
\setlength{\parskip}{4pt}
\setlength{\parindent}{15pt}
\allowdisplaybreaks[2]
\newtheorem{theorem}{Theorem}[section]
\newtheorem{lemma}[theorem]{Lemma}
\newtheorem{corollary}[theorem]{Corollary}
\newtheorem{proposition}[theorem]{Proposition}
\theoremstyle{remark}\newtheorem{remark}[theorem]{Remark}
\newcommand{\E}{\mathbb E}
\newcommand{\Var}{\operatorname{Var}}
\newcommand{\Prob}{\mathbb P}
\newcommand{\dd}{\mathrm d}
\newcommand{\fall}[2]{(#1)_{\underline{#2}}}
\newcommand{\Normal}{\mathcal N}
\newcommand{\coef}[2]{[#1]#2}
\title{Diagonal multiple alignments\\\large All fixed orders, column counts and inverse thresholds}
\author{Report 177 \quad OEIS A262810}
\date{3 October 2026}
\begin{document}
\maketitle
\begin{abstract}
We study binary matrices with $n$ labelled rows, every row sum $n$, ordered nonzero columns, and no prescribed number of columns. Their number is OEIS A262810. We prove a constructive all-fixed-orders expansion, uniformly with every fixed number of real marking derivatives. Centering the roots of a falling factorial explains why only powers of $n^{-2}$ occur after a natural gamma normalization. Three correction polynomials are explicit. The column count has a Gaussian limit on scale $n$ with an essential $n/4$ centering shift. Substituting $k=n$ in the classical fixed-$k$ alignment equivalent misses a factor $e^{-1/12}$. A Lambert-$W$ inverse expansion gives a boundary-safe two-ceiling enclosure. The leading diagonal equivalent is credited to Kot\v{e}\v{s}ovec (2016); the exact alignment family and fixed-dimensional asymptotic are classical. Exact rational finite checks accompany the analytic proof.
\end{abstract}
\begingroup
\small
\setlength{\parskip}{0pt}
\tableofcontents
\endgroup
\clearpage

\section{The objects and the principal results}
For $n\geq1$, let $\mathcal A_n$ be the finite set of binary matrices with $n$ labelled rows, row sums all $n$, and no zero column. Column order matters. If $\ell(M)$ denotes the number of columns, set
\[
 A_n(u)=\sum_{M\in\mathcal A_n}u^{\ell(M)},\qquad a_n=A_n(1),\qquad N=n^2.
\]
Thus $n\leq\ell(M)\leq n^2$. An alignment column records which of $n$ strings advances at that position. The objects are global positional alignments of strings of length $n$; no letter identities, scores, or optimality condition enter the count. The diagonal $a_n$ is A262810, in the array A262809~\cite{oeisdiag,oeisarray}.

Throughout, a fixed order means an integer chosen before $n\to\infty$. All asymptotic constants may depend on that order and on a specified compact parameter interval. No effective onset or numerical remainder constant is asserted.

For $u>0$, put
\begin{equation}\label{eq:carrier}
 \lambda(u)=\log(1+u^{-1}),\qquad
 \mathcal B_n(u)=\frac{\Gamma(N+1)}{(1+u)\Gamma(n+1)^n\lambda(u)^{N+1}}
 \exp\left(-\frac{n-1}{2}\lambda(u)-\frac{\lambda(u)^2}{24}\right).
\end{equation}
\begin{theorem}[Uniform marked expansion]\label{thm:main}
There are rational polynomials $L_j(z)$, constructible at every order, such that for every $J,r\geq0$ and compact $U\subset(0,\infty)$,
\begin{equation}\label{eq:main}
 \left\|\log\frac{A_n(u)}{\mathcal B_n(u)}-
 \sum_{j=1}^J\frac{L_j(\lambda(u))}{N^j}\right\|_{C^r(U)}
 =O_{J,r,U}(N^{-J-1}).
\end{equation}
Here the norm is the maximum of the suprema of derivatives of orders $0,\ldots,r$. In particular,
\begin{align}
 L_1(z)&=\frac{z^4}{2880},\label{eq:L1}\\
 L_2(z)&=\frac{z^4(693-z^2)}{181440},\label{eq:L2}\\
 L_3(z)&=\frac{z^4(3z^4+3200z^2+493920)}{29030400}.\label{eq:L3}
\end{align}
Only integer powers of $N^{-1}=n^{-2}$ occur in this normalization.
\end{theorem}
The statement is an asymptotic family: it does not assert convergence of the infinite series. Section~\ref{sec:proof} supplies the growing-degree lattice and differentiated remainder arguments; Section~\ref{sec:algorithm} makes the coefficient construction explicit.

At $u=1$ write $\lambda=\log 2$ once and for all. The leading equivalent recorded in A262810 by V\'aclav Kot\v{e}\v{s}ovec on 23 March 2016 is
\begin{equation}\label{eq:K}
 K_n=\frac{n^{n^2-n/2+1}}
 {e^{1/12}\,2^{n+\lambda/24}\,\pi^{(n-1)/2}\lambda^{n^2+1}}.
\end{equation}
Its factors $e^{-1/12}$ and $e^{-\lambda^2/24}$ are part of that credited equivalent.
\begin{corollary}[Corrections to the credited equivalent]\label{cor:C}
For every fixed $J$, $\log(a_n/K_n)=\sum_{j=1}^J C_jn^{-2j}+O(n^{-2J-2})$. The first three logarithmic coefficients are
\begin{align}
 C_1&=\frac{\lambda^4+248}{2880},\label{eq:C1}\\
 C_2&=-\frac{\lambda^6-693\lambda^4+144}{181440},\label{eq:C2}\\
 C_3&=\frac{3\lambda^8+3200\lambda^6+493920\lambda^4-63360}{29030400}.\label{eq:C3}
\end{align}
The first three ordinary multiplicative coefficients are respectively
$C_1$, $C_2+C_1^2/2$, and $C_3+C_1C_2+C_1^3/6$.
\end{corollary}

Let $\mathsf L_n$ be the number of columns under the uniform law on $\mathcal A_n$.
\begin{theorem}[Column count]\label{thm:law}
Set $\sigma^2=(1-\lambda)/(4\lambda^2)>0$. Then
\begin{align}
 \E\mathsf L_n={}&\frac{n^2}{2\lambda}+\frac n4+
 \left(\frac1{2\lambda}-\frac34+\frac\lambda{24}\right)
 -\frac{\lambda^3}{1440n^2}+O(n^{-4}),\label{eq:mean}\\
 \Var(\mathsf L_n)={}&\sigma^2 n^2-\frac n8+
 \left(\sigma^2-\frac{7+\lambda}{48}\right)
 +\frac{3\lambda^2+\lambda^3}{2880n^2}+O(n^{-4}),\label{eq:variance}\\
 \frac{\mathsf L_n-n^2/(2\lambda)-n/4}{n}
 &\ \Longrightarrow\ \Normal(0,\sigma^2).\label{eq:clt}
\end{align}
All fixed cumulants have constructive expansions by differentiating \eqref{eq:main}. Centering only by $n^2/(2\lambda)$ changes the limiting mean in \eqref{eq:clt} to $1/4$.
\end{theorem}

\section{Exact constructions and finite certificates}\label{sec:exact}
For arbitrary row sums $n_1,\ldots,n_k$, the generating function of a sequence of nonzero binary columns is
\begin{equation}\label{eq:gf}
 \sum_{n_1,\ldots,n_k\geq0} A(n_1,\ldots,n_k;u)
 x_1^{n_1}\cdots x_k^{n_k}
 =\frac1{1-u(\prod_{i=1}^k(1+x_i)-1)}.
\end{equation}
For $u>0$, expanding near $x=0$ gives the positive representation
\begin{equation}\label{eq:positive}
 A(n_1,\ldots,n_k;u)=\frac1{1+u}
 \sum_{m\geq0}\left(\frac u{1+u}\right)^m
 \prod_{i=1}^k\binom m{n_i}.
\end{equation}
The unmarked identity is classical; Eger~\cite[Section 5.5]{eger} gives both this and inclusion--exclusion. It is also recorded for A262809~\cite{oeisarray}. The marking follows the same column construction.

Let $p_n(m)=\binom mn^n$. Inclusion--exclusion over columns unused by every row gives
\begin{equation}\label{eq:difference}
 d_{n,j}=[u^j]A_n(u)=\Delta^jp_n(0)
 =\sum_{i=0}^j(-1)^{j-i}\binom ji\binom in^n,
 \qquad A_n(u)=\sum_{j=0}^{N}d_{n,j}u^j.
\end{equation}
The signed formula is evaluated exactly with integers, not by cancellation of floating-point approximations. A genuinely positive independent construction expands $p_n$ in the binomial basis using
\begin{equation}\label{eq:basis}
 \binom mr\binom mn
 =\sum_{j=\max(r,n)}^{r+n}\binom jr\binom r{r+n-j}\binom mj.
\end{equation}
Indeed, choose an $r$-set and an $n$-set in an $m$-set by first choosing their $j$-element union, then the $r$-set, then the $r+n-j$ intersection. Repeating \eqref{eq:basis} $n$ times yields the same $d_{n,j}$ with positive integer operations.

A third, small-instance recursion is
\begin{equation}\label{eq:grid}
 F(v)=\sum_{\varnothing\ne S\subseteq\operatorname{supp}(v)}F(v-\boldsymbol1_S),
 \qquad F(0)=1,
\end{equation}
where $v$ is the vector of remaining row sums. The first nonzero column determines the summand, proving the recursion.

For a rational mark $u>0$, write $\rho=u/(1+u)$ and
\[
 t_m=\frac{\rho^m}{1+u}\binom mn^n,\qquad
 \frac{t_{m+1}}{t_m}=\rho\left(\frac{m+1}{m+1-n}\right)^n=:r_m
 \quad(m\geq n).
\]
The ratios decrease to $\rho<1$. If $r_M<1$ and $S_M=\sum_{m=n}^M t_m$, then
\begin{equation}\label{eq:tail}
 S_M\leq A_n(u)\leq S_M+\frac{t_Mr_M}{1-r_M}.
\end{equation}
Every endpoint is rational. These finite certificates check exact identities; they are logically separate from the analytic asymptotic theorem.

\section{Centering the growing degree}\label{sec:center}
Put
\[
 h=\frac{n-1}2,\quad q=m-h,\quad a_j=h-j\ (0\leq j<n),\quad
 P_n(q)=\prod_{j=0}^{n-1}(q+a_j)^n.
\]
The multiset of roots is symmetric. For $q>h$,
\begin{equation}\label{eq:rootexp}
 P_n(q)=q^N\exp\left(-\sum_{r\geq1}B_r(N)q^{-2r}\right),\qquad
 B_r(N)=\frac n{2r}\sum_{j=0}^{n-1}(h-j)^{2r}.
\end{equation}
Also $0\leq P_n(q)\leq q^N$ for $q\geq h$, by pairing $q+a$ with $q-a$. At $q=h$, the polynomial has a zero of multiplicity $n$.

Let $\mathrm B_s(x)$ be the Bernoulli polynomial, distinguished typographically from $B_r(N)$. The power sum is
\[
 \sum_{j=0}^{n-1}(h-j)^{2r}
 =\frac{\mathrm B_{2r+1}((n+1)/2)-\mathrm B_{2r+1}((1-n)/2)}{2r+1}.
\]
It is an odd polynomial in $n$, of degree $2r+1$, as follows either by replacing $n$ with $-n$ or by Bernoulli reflection. Thus $B_r(N)$ is a rational polynomial of degree $r+1$ in $N$. In particular,
\begin{align}
 B_1(N)&=N(N-1)/24,\notag\\
 B_2(N)&=N(N-1)(3N-7)/960,\notag\\
 B_3(N)&=N(N-1)(3N^2-18N+31)/8064,\label{eq:B}\\
 B_4(N)&=N(N-1)(5N^3-55N^2+239N-381)/92160.\notag
\end{align}
This parity is the source of the all-orders $n^{-2}$ structure; observing the first coefficients alone would not prove it.

\section{Uniform analytic remainder bounds}\label{sec:proof}
Fix $0<\alpha\leq\lambda\leq\beta<\infty$. Constants below are uniform on this interval. Define
\[
 f(q)=e^{-\lambda q}P_n(q)\boldsymbol1_{q\geq h},\qquad
 G_N(\lambda)=\frac{\Gamma(N+1)}{\lambda^{N+1}},\qquad
 I_n(\lambda)=\int_h^\infty e^{-\lambda q}P_n(q)\,\dd q.
\]
The lattice is $\mathbb Z-h$. Its point $h$ has value zero; all original terms with $m<n$ vanish. Equation~\eqref{eq:positive} becomes
\begin{equation}\label{eq:lattice}
 A_n(u)=\frac{e^{-\lambda h}}{(1+u)(n!)^n}
 \sum_{q\in\mathbb Z-h}f(q).
\end{equation}

\subsection{Lattice replacement with growing degree}
\begin{lemma}\label{lem:em}
For all fixed $M,r\geq0$,
\[
 \left\|\frac{\sum_{q\in\mathbb Z-h}f(q)-I_n(\lambda)}{G_N(\lambda)}
 \right\|_{C^r([\alpha,\beta])}=O_{M,r}(N^{-M}).
\]
Moreover $cG_N\leq I_n\leq G_N$ with $c>0$ independent of $n$ and $\lambda$ for large $n$.
\end{lemma}
\begin{proof}
For any fixed integer $R\geq2$ and $n>R$, the cutoff at $h$ has the boundary regularity required by Euler--Maclaurin. Its derivatives through order $R-1$ vanish there, and its $R$th derivative is integrable. The periodic-Bernoulli remainder, or $R$ integrations by parts in Poisson summation, gives
\begin{equation}\label{eq:em}
 \left|\sum_{q\in\mathbb Z-h}f(q)-I_n\right|
 \leq c_R\|f^{(R)}\|_1.
\end{equation}
We bound the integrated derivative in three regions; derivatives here are with respect to $q$.

On $\mathcal C=\{|q-N/\lambda|\leq\sqrt N\log N\}$, all $q+a_j$ are comparable to $N$. With $\ell=\log f$,
\[
 \ell'=-\lambda+n\sum_j(q+a_j)^{-1}
 =O(N^{-1/2}\log N),\qquad
 \ell^{(s)}=O_s(N^{1-s})\quad(s\geq2).
\]
For the first equality, the linear $a_j$ term cancels; the correction to $N/q$ is $O(N^2/q^3)=O(N^{-1})$. In Bell's derivative formula, a product with $m_s$ factors $\ell^{(s)}$ and $\sum s m_s=R$ is at most $O_R(N^{-R/2}(\log N)^R)$. Hence
\begin{equation}\label{eq:centralderivative}
 |f^{(R)}|\leq c_R N^{-R/2}(\log N)^R f\quad\hbox{on }\mathcal C.
\end{equation}
The logarithm of $P_n(q)/q^N$ is bounded between fixed constants there by \eqref{eq:rootexp}. A gamma variable of shape $N+1$ and rate $\lambda$ lies in $\mathcal C$ with probability tending to one uniformly. This proves $I_n\geq cG_N$; the upper bound follows from $P_n\leq q^N$.

On $[N/(2\lambda),2N/\lambda]\setminus\mathcal C$, the same logarithmic derivatives are bounded by fixed powers of $N$. The gamma Chernoff bound gives a loss $\exp(-c\log^2N)$ after integration, smaller than $N^{-A}$ for every fixed $A$. For completeness, the bound follows from the gamma mgf $(1-t/\lambda)^{-(N+1)}$ by exponential Markov inequality, optimizing $t$; on compact relative deviations the rate is quadratic.

On the two far tails, do not use singular logarithmic derivatives near $h$. Differentiating the product of its $N$ nonnegative linear factors yields
\begin{equation}\label{eq:leibniz}
 |f^{(R)}(q)|\leq\sum_{j=0}^R\binom Rj\lambda^{R-j}
 \fall Nj e^{-\lambda q}(q+h)^{N-j}.
\end{equation}
Integrate each term after the substitution $z=q+h$. If $Z_j$ has gamma shape $N-j+1$ and rate $\lambda$, after factoring out $\binom Rj$, the corresponding integral divided by $G_N$ is exactly
\begin{equation}\label{eq:tailintegrated}
 \lambda^R e^{\lambda h}\Prob(Z_j\hbox{ belongs to the shifted integration region}),
\end{equation}
since $\fall Nj\Gamma(N-j+1)/\Gamma(N+1)=1$. In the lower tail the upper endpoint is $N/(2\lambda)+h$; in the upper tail the lower endpoint is $2N/\lambda+h$. For fixed $j$ these probabilities are $O(e^{-cN})$. The prefactor $e^{\lambda h}=e^{O(\sqrt N)}$ is absorbed. This integrated estimate avoids a false pointwise comparison at the cutoff.

Combining the regions proves
\[
 \frac{\|f^{(R)}\|_1}{G_N}
 \leq c_R N^{-R/2}(\log N)^R+O_A(N^{-A}).
\]
Choose $R>2M$ and then slightly larger to absorb logarithms. For $r$ parameter derivatives, differentiation inserts $(-q)^r$; Leibniz bounds for the $q$-derivatives of $q^rf$ lose at most fixed powers $N^r$. Derivatives of $G_N^{-1}$ likewise lose only fixed powers of $N$. Increase $R$ again. The zero at the boundary retains sufficient order for each fixed $r$, completing the $C^r$ assertion.
\end{proof}

\subsection{The gamma expectation and its differentiated expansion}
Set $\varepsilon=N^{-1}$ and let $Y$ have gamma shape $N+1$ and rate $N$, so its law is independent of $\lambda$. On $y>\lambda h/N$ put
\[
 Q_N(y,\lambda)=\prod_j\left(1+\frac{\lambda a_j}{Ny}\right)^n,
\]
and put it equal to zero below that cutoff. Scaling $q=Ny/\lambda$ gives
\begin{equation}\label{eq:expect}
 I_n/G_N=\E Q_N(Y,\lambda).
\end{equation}
Fix a neighborhood $V$ of $1$ with compact closure in $(0,\infty)$. On $V$, the series \eqref{eq:rootexp} is absolutely convergent and
\begin{equation}\label{eq:g}
 \log Q_N(y,\lambda)=-\frac a{y^2}+\sum_{j=1}^J\varepsilon^j g_j(y;\lambda)+O(\varepsilon^{J+1}),
 \qquad a=\lambda^2/24.
\end{equation}
The expansion through $\varepsilon^J$ has remainder $O(\varepsilon^{J+1})$ in every fixed mixed $y,\lambda$ derivative norm. Indeed,
\[
 B_r(N)(Ny/\lambda)^{-2r}\leq \frac N{2r}
 \left(\frac{\lambda h}{Ny}\right)^{2r},
\]
so the tail starting at $r=J+2$ is $O(N^{-J-1})$. Derivatives contribute only fixed powers of $r$ and bounded reciprocal powers of $y,\lambda$, absorbed by the geometric tail. Every retained $B_r$ is a polynomial in $N$.

The expectation outside $V$ is negligible in every required fixed $\lambda$ derivative norm. Here is an integrated bound that also handles the moving cutoff. For a derivative order $s$, the product rule and $|a_j|\leq h$ give, on the positive domain,
\[
 |\partial_\lambda^s Q_N(y,\lambda)|
 \leq\fall Ns\left(\frac h{Ny}\right)^s
 \left(1+\frac{\beta h}{Ny}\right)^{N-s}.
\]
Multiplication by the gamma density and the substitution $z=y+\beta h/N$ bound each tail integral by
\begin{equation}\label{eq:scaledtail}
 h^s e^{\beta h}\Prob\{\operatorname{Gamma}(N-s+1,N)
 \hbox{ lies in the shifted tail}\}.
\end{equation}
The factorials cancel exactly. Because the shift is $O(N^{-1/2})$ and $V$ is fixed, these probabilities are $O(e^{-cN})$. Boundary terms from differentiating the cutoff vanish for $n>s$, owing to the multiplicity-$n$ zero. This proves the differentiated tail assertion without assuming complex-parameter uniformity.

Exponentiate \eqref{eq:g} as
\[
 Q_N(y,\lambda)=e^{-a/y^2}\sum_{q=0}^JH_q(y;\lambda)\varepsilon^q
 +O_{C^r}(\varepsilon^{J+1})\quad(y\in V).
\]
For the $q$th term put $p=J-q$. Taylor-expand its smooth coefficient about $y=1$ through degree $2p+1$. The next, even-degree absolute remainder has expectation
\[
 O\bigl(\E|Y-1|^{2p+2}\bigr)=O(\varepsilon^{p+1}).
\]
This step needs the odd degree $2p+1$ before estimating the remainder; stopping at degree $2p$ with an absolute odd moment would not suffice for the stated order. The signed moments are exactly
\begin{equation}\label{eq:mu}
 \mu_r(\varepsilon)=\E(Y-1)^r
 =\sum_{j=0}^r(-1)^{r-j}\binom rj\prod_{i=1}^j(1+i\varepsilon).
\end{equation}
Their valuation is at least $\lceil r/2\rceil$. For example, $Y-1$ has cumulants $\kappa_1=\varepsilon$ and
$\kappa_r=(r-1)!(\varepsilon^{r-1}+\varepsilon^r)$ for $r\geq2$; the moment--cumulant formula proves the valuation and the even absolute-moment estimate. Thus the degree $2p+1$ term can be omitted from the coefficient algorithm, because it is already $O(\varepsilon^{p+1})$.

All these estimates commute with any fixed number of $\lambda$ derivatives: the law of $Y$ is independent of $\lambda$, and coefficient bounds are uniform. Taking logarithms is valid because the leading expectation $e^{-a}$ stays bounded away from zero. Lemma~\ref{lem:em}, \eqref{eq:lattice}, and the smooth change of variable $\lambda=\log(1+1/u)$ prove Theorem~\ref{thm:main}.

\section{A constructive algorithm at every fixed order}\label{sec:algorithm}
The following finite rational procedure computes $L_1,\ldots,L_J$.
\begin{enumerate}[label=\arabic*.]
\item Compute $B_r(N)$ for $1\leq r\leq J+1$ by the Bernoulli formula in Section~\ref{sec:center}. Expand
\[
-\sum_{r=1}^{J+1}B_r(\varepsilon^{-1})\lambda^{2r}\varepsilon^{2r}y^{-2r}
\]
through degree $J$, giving $-a/y^2+\sum_{j=1}^Jg_j\varepsilon^j$.
\item Form $H_0=1$ and
\begin{equation}\label{eq:H}
 H_j=\frac1j\sum_{s=1}^j s g_s H_{j-s}.
\end{equation}
\item Compute the moments \eqref{eq:mu} through degree $2J$ and form, discarding powers above $\varepsilon^J$,
\begin{equation}\label{eq:algorithm}
 \sum_{q=0}^J\varepsilon^q\sum_{r=0}^{2(J-q)}\frac{\mu_r(\varepsilon)}{r!}
 \left[e^a\frac{\dd^r}{\dd y^r}
 \{e^{-a/y^2}H_q(y;\lambda)\}\right]_{y=1}
 =\sum_{j=0}^JR_j(\lambda)\varepsilon^j.
\end{equation}
No transcendental symbolic operation is needed: start from $f=H_q$ and replace
$f$ by $\partial_y f+2ay^{-3}f$ at each differentiation, then evaluate at $y=1$.
\item With $R_0=1$, extract the logarithm by
\begin{equation}\label{eq:logrec}
 L_j=R_j-\frac1j\sum_{s=1}^{j-1}s L_sR_{j-s}.
\end{equation}
\end{enumerate}
Every operation is rational polynomial or Laurent-polynomial arithmetic. The proof in Section~\ref{sec:proof} establishes the remainder at each fixed order, rather than merely identifying a formal series.

As a transparent first-order check,
\[
 g_1=ay^{-2}-\frac{\lambda^4}{320}y^{-4}.
\]
For $f(y)=e^{-a/y^2}$, the gamma expectation correction is
$f'(1)+f''(1)/2=e^{-a}(2a^2-a)$. Adding $g_1(1)$ cancels the terms in $a$ and leaves
\[
 L_1=2a^2-\lambda^4/320=\lambda^4/2880.
\]
The supplied rational verifier implements this algorithm and checks the printed $L_j$.

To pass from the exact gamma carrier to \eqref{eq:K}, set
$b_s=\mathrm B_{2s}/(2s(2s-1))$, with $\mathrm B_{2s}$ the Bernoulli number. Stirling's expansion in
$\log\Gamma(N+1)-n\log\Gamma(n+1)$ gives
\begin{equation}\label{eq:stirling}
 C_j=L_j(\lambda)-b_{j+1}+\mathbf{1}_{\{j\text{ odd}\}}b_{(j+1)/2}.
\end{equation}
The constant from $-n\log\Gamma(n+1)$ is $-1/12$. The remaining leading terms give exactly \eqref{eq:K}, including $e^{-\lambda^2/24}$. Taking sufficiently many Stirling terms proves the remainder in Corollary~\ref{cor:C}; \eqref{eq:stirling} gives \eqref{eq:C1}--\eqref{eq:C3}.

\section{The complete marked law and its central limit}\label{sec:statistics}
The polynomial \eqref{eq:difference} is the exact column-count law:
\[
 \Prob(\mathsf L_n=j)=d_{n,j}/a_n,\qquad
 \E e^{t\mathsf L_n}=A_n(e^t)/A_n(1).
\]
For all real $t$ in any fixed compact interval, write $\lambda(t)=\log(1+e^{-t})$. Taking logarithms in Theorem~\ref{thm:main} gives the complete all-orders real cumulant-generating expansion
\begin{align}\label{eq:cgf}
 \log\E e^{t\mathsf L_n}
 ={}&-(N+1)\log\frac{\lambda(t)}\lambda
 -\log\frac{1+e^t}{2}
 -h(\lambda(t)-\lambda)
 -\frac{\lambda(t)^2-\lambda^2}{24}\notag\\
 &+\sum_{j=1}^J\frac{L_j(\lambda(t))-L_j(\lambda)}{N^j}
 +O_{C^r}(N^{-J-1}).
\end{align}
The remainder vanishes at $t=0$ and has the proved fixed derivative bounds. For a fixed positive mark $u_0$, the tilted law assigning probability $u_0^{\ell(M)}/A_n(u_0)$ has the same formula with $A_n(u_0e^t)/A_n(u_0)$. Thus \eqref{eq:main} describes all fixed positive tilts as well as all fixed cumulants.

\subsection{Direct moment identities and cancellation control}
Introduce an auxiliary integer $\mathsf M$ with probability proportional to $2^{-m}\binom mn^n$. This variable is a device from the positive sum and is not the column count. Differentiating the absolutely convergent exact series in \eqref{eq:positive}, then setting $u=e^t$ and $t=0$, proves
\begin{equation}\label{eq:Mrelations}
 \E\mathsf L_n=\frac{\E\mathsf M-1}{2},\qquad
 \Var(\mathsf L_n)=\frac{\Var(\mathsf M)-\E\mathsf M-1}{4}.
\end{equation}
There is also an exact probabilistic explanation. Sample $\mathsf M=m$ with the stated mass, choose each row's $n$ positions independently and uniformly among $m$ columns, then delete the zero columns. For a specified length-$\ell$ alignment, the total embedding weight is
\[
 \sum_{m\geq\ell}2^{-m}\binom m\ell=2,
\]
independent of the alignment, so deletion gives the uniform law on $\mathcal A_n$. Conditional on $\mathsf L_n=\ell$, the difference $\mathsf M-\ell$ is the sum of $\ell+1$ independent geometric failure counts with success probability $1/2$. Therefore
\[
 \E(\mathsf M\mid\mathsf L_n)=2\mathsf L_n+1,\qquad
 \Var(\mathsf M\mid\mathsf L_n)=2(\mathsf L_n+1).
\]
Total expectation and total variance recover \eqref{eq:Mrelations}.

For clarity, $\lambda'(0)=-1/2$ and $\lambda''(0)=1/4$; the separate term $-\log(1+e^t)$ contributes $-1/2$ and $-1/4$ respectively.

Put $q=\mathsf M-h$. The $C^2$ version of the theorem justifies differentiating the normalized logarithm of its positive lattice partition function. It gives
\begin{align}
 \E q&=\frac{N+1}{\lambda}+\frac\lambda{12}
 -\frac{\lambda^3}{720N}+O(N^{-2}),\label{eq:qmean}\\
 \Var(q)&=\frac{N+1}{\lambda^2}-\frac1{12}
 +\frac{\lambda^2}{240N}+O(N^{-2}).\label{eq:qvar}
\end{align}
These formulas can also be obtained without differentiating a remainder: insert $m$ and $m^2$ separately in the positive series, apply the same lattice estimate, then use $m=h+NY/\lambda$ in the gamma integral. The polynomial insertions preserve the Taylor proof. One must retain enough terms before subtracting the squared mean: an $m^2$ insertion has scale $N^2$, so to get an $O(N^{-2})$ absolute remainder it suffices to keep the normalized expectation through $\varepsilon^3$ with remainder $O(\varepsilon^4)$, and correspondingly enough terms in the first moment and denominator. The arbitrary-order theorem allows this precision. A bare $O(N^{-2})$ normalized second-moment error would not justify the variance claim.

Now substitute $\E\mathsf M=\E q+h$ and $\Var(\mathsf M)=\Var(q)$ in \eqref{eq:Mrelations}. This proves \eqref{eq:mean}--\eqref{eq:variance}, including the constant terms and the $n^{-2}$ corrections.

\subsection{A real moment-generating-function proof}
For a fixed real $s$, take $t=s/n$ in \eqref{eq:cgf}. The function
$g(t)=-\log(\lambda(t)/\lambda)$ satisfies
\[
 g'(0)=\frac1{2\lambda},\qquad
 g''(0)=\frac{1-\lambda}{4\lambda^2}=\sigma^2.
\]
Hence $Ng(s/n)=sn/(2\lambda)+\sigma^2s^2/2+O(n^{-1})$. The $h$-term is
\[
 -h(\lambda(s/n)-\lambda)=s/4+O(n^{-1}).
\]
All other differences in \eqref{eq:cgf} tend to zero. Therefore
\[
 \log\E\exp\left\{s\frac{\mathsf L_n-N/(2\lambda)-n/4}{n}\right\}
 \longrightarrow\frac{\sigma^2s^2}{2}.
\]
The real mgf convergence theorem proves \eqref{eq:clt}. This argument requires no complex extension of the marking theorem. It provides neither a local central limit theorem nor a Berry--Esseen rate, and no such conclusion is claimed.

\section{The fixed dimensional limit and the diagonal}\label{sec:fixed}
For fixed $k$, let $A(n,k)=A(n,\ldots,n;1)$ with $k$ row sums. The classical theorem of Griggs, Hanlon, Odlyzko, and Waterman~\cite{ghow}, as displayed also by Pemantle and Wilson~\cite[Section 4.9]{pw} and Eger~\cite{eger}, is
\begin{equation}\label{eq:G}
 A(n,k)\sim G(n,k),\qquad
 G(n,k)=\frac{2^{(1-k^2)/(2k)}(2^{1/k}-1)^{-kn-1}}
 {\sqrt k\,(\pi n)^{(k-1)/2}}\quad(n\to\infty,\ k\hbox{ fixed}).
\end{equation}
Its quantifier on $k$ matters.
\begin{proposition}\label{prop:fixed}
Substitution $k=n$ in the fixed-dimensional leading expression gives
\[
 \log\frac{a_n}{G(n,n)}=-\frac1{12}
 +\frac{31/360+\lambda^2/24}{n^2}+O(n^{-4}),
 \qquad \frac{a_n}{G(n,n)}\longrightarrow e^{-1/12}.
\]
\end{proposition}
\begin{proof}
Use
\[
 \log\frac{e^z-1}{z}=\frac z2+\frac{z^2}{24}
 -\frac{z^4}{2880}+O(z^6),\qquad z=\lambda/n,
\]
in the exact expression $G(n,n)$, and compare with \eqref{eq:K} and $C_1$ in \eqref{eq:C1}. The constant discrepancy is $-1/12$; the $n^{-2}$ coefficient is $31/360+\lambda^2/24$ after the $\lambda^4$ terms cancel.
\end{proof}
This is a nonuniformity example, not an error in the fixed-$k$ theorem. The previously reported typographical corrections to a Waterman handbook formula discussed in~\cite{pw} are a separate matter and are not claimed here.

\section{Lambert inversion with a safe integer boundary}\label{sec:inverse}
For $\lambda=\log2$, define
\begin{align}
 F_0(x)&=x^2\log(x/\lambda),\notag\\
 H_0(x)&=-\frac x2\log x-x\log2-\frac{x-1}{2}\log\pi
 +\log x-\log\lambda-\frac1{12}-\frac{\lambda^2}{24},\label{eq:FH}\\
 H_J(x)&=H_0(x)+\sum_{j=1}^JC_jx^{-2j},\qquad F_J(x)=F_0(x)+H_J(x).\notag
\end{align}
Corollary~\ref{cor:C} says $\log a_n=F_J(n)+O(n^{-2J-2})$. These are smooth asymptotic surrogates, not a claimed analytic continuation of the discrete sequence.

\subsection{All fixed inverse orders}
The positive Lambert branch solves $W(z)e^{W(z)}=z$ for $z>0$. The large inverse of $F_0$ is
\begin{equation}\label{eq:x0}
 x_0(T)=\sqrt{\frac{2T}{W(2T/\lambda^2)}}.
\end{equation}
For large $T$, let $x_J(T)$ be the exact large real root of $F_J(x)=T$ and set
$D=(F_0'(x))^{-1}\dd/\dd x$.
\begin{theorem}[Constructive inverse expansion]\label{thm:inverse}
For every fixed $J\geq0$ and $K\geq1$,
\begin{equation}\label{eq:inverse}
 x_J(T)=x_0(T)+\sum_{k=1}^K\frac{(-1)^k}{k!}
 \left.D^{k-1}\left(\frac{H_J(x)^k}{F_0'(x)}\right)\right|_{x=x_0(T)}
 +O_{J,K}(x_0(T)^{-K}).
\end{equation}
In particular, with the functions on the right evaluated at $x_0$,
\begin{equation}\label{eq:inverse2}
 x_J=x_0-\frac{H_J}{F_0'}+
 \frac{H_JH_J'}{(F_0')^2}-\frac{H_J^2F_0''}{2(F_0')^3}
 +O(x_0^{-2}).
\end{equation}
The first correction tends to $1/4$, with logarithmic corrections.
\end{theorem}
\begin{proof}
For $0\leq\rho\leq1$ solve $F_0(X)+\rho H_J(X)=T$ near $x_0$. Since $H_J/F_0'=O(1)$ and $F_0'\sim2x\log x$, monotonicity and the intermediate value theorem give a unique root $X=x_0+O(1)$ uniformly in $\rho$. The derivative $F_0'+\rho H_J'$ remains comparable to $F_0'$.

Uniformly in this neighborhood,
\[
 \frac{H_J^{(\ell)}}{F_0'}=O(x_0^{-\ell})\quad(\ell\geq1),\qquad
 \frac{F_0^{(\ell)}}{F_0'}=O(x_0^{1-\ell})\quad(\ell\geq2).
\]
Implicit differentiation, followed by induction, gives
$\partial_\rho^kX=O_k(x_0^{1-k})$ for every fixed $k\geq1$. In the differentiated equation the term containing the highest derivative is $(F_0'+\rho H_J')\partial_\rho^kX$; the other Bell-polynomial terms obey the displayed bounds and already established derivative bounds. Taylor's theorem in $\rho$ through degree $K$ consequently has remainder $O(x_0^{-K})$.

To compute its coefficients, make the change $w=F_0(x)$. The equation becomes $w=T-\rho H_J(F_0^{-1}(w))$. The Lagrange inversion formula applied to $F_0^{-1}(w)$ gives exactly the finite coefficients in \eqref{eq:inverse}; they can equivalently be obtained by successive implicit differentiation at $\rho=0$. The first two yield \eqref{eq:inverse2}. Finally $-H_J/F_0'\to1/4$.
\end{proof}
No convergence of the infinite inverse series follows or is needed.

\subsection{The integer threshold and near-boundary ambiguity}
Define
\[
 \nu(T)=\min\{n\geq1:a_n\geq e^T\}.
\]
The sequence is strictly increasing for positive $n$. To see this, add a zero row to an $n$-row matrix, append a column with ones in all old rows and zero in the new row, then append $n+1$ columns supported only on the new row. Every row sum becomes $n+1$. Deleting the fixed final $n+2$ columns and the new row recovers the original. This is an injection; the all-ones $(n+1)\times(n+1)$ matrix is outside its image, proving strictness.

\begin{theorem}[Boundary-safe enclosure]\label{thm:ceiling}
For every fixed $J\geq0$ there exist constants $C_J>0$ and $T_J$ such that, with $x=x_J(T)$ and
\[
 \delta_J(x)=C_J\frac{x^{-2J-3}}{\log x},
\]
all $T\geq T_J$ satisfy
\begin{equation}\label{eq:ceiling}
 \left\lceil x-\delta_J(x)\right\rceil\leq\nu(T)
 \leq\left\lceil x+\delta_J(x)\right\rceil.
\end{equation}
For large $T$ the two integer endpoints differ by at most one. At a sequence threshold $T=\log a_n$,
$x_J(T)-n=O(n^{-2J-3}/\log n)$.
\end{theorem}
\begin{proof}
Write $\log a_m=F_J(m)+r_m$ with $|r_m|\leq C m^{-p}$, $p=2J+2$, for all large integers $m$. Since $F_J'\asymp x\log x$ near $x$ and $F_J(x)=T$, a sufficiently large fixed $C_J$ gives, for integers $|m-x|\leq2$,
\[
 m<x-\delta_J(x)\ \Longrightarrow\ \log a_m<T,
 \qquad
 m\geq x+\delta_J(x)\ \Longrightarrow\ \log a_m\geq T.
\]
First applying the same derivative bounds at neighboring integers a fixed distance from $x$ locates the threshold in this neighborhood. Monotonicity then extends the implications to more distant integers, proving \eqref{eq:ceiling}. The threshold-sample assertion follows by the mean value theorem.
\end{proof}
If the two ceilings coincide, they determine the exact threshold. A universal single-ceiling formula is unwarranted when $T$ is arbitrarily close to $\log a_n$. To use the $K$-term approximation \eqref{eq:inverse} instead of the exact surrogate root, enlarge the radius by its error. Taking $K\geq2J+4$ makes $O(x_0^{-K})$ smaller than the $x^{-2J-3}/\log x$ scale. Taking only $K=2J+3$ does not establish this extra logarithmic gain. The constants and onsets in this enclosure are existence statements, not certified numerical thresholds.

\section{A binomial coupling for the nonnegative analogue}\label{sec:corollary}
Let $\widetilde A_n(u)$ count nonnegative-integer matrices with the same labelled rows and row sums, ordered columns, and no zero column. This is the diagonal A316677 of the array A316674. The already recorded unmarked relation to binary matrices is
$\widetilde A_n(1)=2^{n-1}a_n$~\cite{oeisnonnegative}. We use it as a corollary of the marked exact construction, not as an independent discovery.

The column generating function replaces $\prod(1+x_i)$ by $\prod(1-x_i)^{-1}$. Hence
\[
 \widetilde A_n(u)=\frac1{1+u}\sum_{m\geq0}\left(\frac u{1+u}\right)^m
 \binom{m+n-1}{n}^{\!n}.
\]
Shifting $k=m+n-1$ gives the exact polynomial identity, interpreted for $u>0$,
\begin{equation}\label{eq:tilde}
 \widetilde A_n(u)=\left(\frac{1+u}{u}\right)^{n-1}A_n(u).
\end{equation}
The apparent negative powers cause no problem because $A_n(u)$ is divisible by $u^n$. Equivalently, under the uniform laws,
\begin{equation}\label{eq:coupling}
 \widetilde{\mathsf L}_n\ \overset{d}=\
 \mathsf L_n-(n-1)+\operatorname{Bin}(n-1,1/2),
\end{equation}
with the binomial variable independent of $\mathsf L_n$. Thus
\[
 \E\widetilde{\mathsf L}_n=\E\mathsf L_n-(n-1)/2,\qquad
 \Var(\widetilde{\mathsf L}_n)=\Var(\mathsf L_n)+(n-1)/4,
\]
and
\[
 \frac{\widetilde{\mathsf L}_n-n^2/(2\lambda)+n/4}{n}
 \Longrightarrow\Normal(0,\sigma^2).
\]
All fixed-order marked expansions follow by adding $(n-1)\log((1+u)/u)$ to the binary logarithm. Likewise the logarithmic count surrogate is $F_J(x)+(x-1)\log2$. The inverse proof applies after replacing $H_J$ by $H_J+(x-1)\log2$, with the same remainder scales and boundary-safe enclosure.

\section{Sources, reproducibility and remaining questions}\label{sec:scope}
\subsection{Attribution and the scope of the source check}
The leading diagonal equivalent \eqref{eq:K} is recorded in OEIS with Kot\v{e}\v{s}ovec's 2016 attribution. The exact family and fixed-dimensional asymptotic precede this report: Griggs--Hanlon--Odlyzko--Waterman~\cite{ghow}, Pemantle--Wilson~\cite{pw}, and Eger~\cite{eger} are the primary sources used. The report's work is the proved growing-diagonal expansion, its explicit corrections, the marked law and shifted Gaussian limit, the fixed-dimension substitution comparison, and the inverse construction.

The inspected OEIS entries and primary papers do not supply this collection of diagonal results. This is a bounded source assessment, not a theorem of global priority or an exhaustive literature search. A262810 cross-references both A262809 and A316677. The directly accessible A316674 page establishes the nonnegative array relation; a direct fetch of A316677 failed. The A262810 b-file also could not be retrieved, so no b-file comparison is claimed. The separate source audit records these limits and the document URLs.

\subsection{What the package verifies}
The companion package contains this complete \LaTeX{} source, the PDF, source metadata, and a Python standard-library verifier using exact integers and fractions. Its independent finite constructions are the positive binomial-basis multiplication \eqref{eq:basis} and the signed difference table \eqref{eq:difference}; the small grid recursion \eqref{eq:grid} supplies a third check. For $1\leq n\leq10$ and $u=1/2,1,2$, the exact rational intervals \eqref{eq:tail} enclose the finite marked polynomial, with relative tail less than $10^{-70}$. The coefficient algorithm checks $B_1$--$B_4$, $L_1$--$L_3$, and $C_1$--$C_3$. An independent Gaussian Laplace calculation agrees with the gamma-central-moment route; the extension run checks both routes through order five. The nonnegative coupling and exact moment relations are checked separately.

The build uses an installed TeX stack, shell escape disabled, isolated temporary directories, fixed timestamps and a complete strict file manifest. Verification and negative guards run with ordinary and optimized Python. A fresh extraction can replay the verification and reproduce the artifacts under the same toolchain. No downloaded program is executed. These are arithmetic and packaging certificates; they do not replace the proof of the analytic remainders or supply effective asymptotic constants.

\subsection{Further questions}
Several extensions remain open within this report's scope.
\begin{enumerate}[label=\arabic*.]
\item Determine the largest growth regimes $k=k(n)$ for which the fixed-dimensional expression \eqref{eq:G} remains relatively accurate, and describe the crossover to $k\asymp n$. The diagonal factor already shows that uniformity cannot extend unchanged to $k=n$.
\item Prove a local central limit theorem and an explicit convergence rate for the column count. The present real-mgf argument alone does not establish either.
\item Extract effective constants and onsets for the all-orders remainder and the inverse enclosure. The rational finite checks do not make the asymptotic constants effective.
\item Study the large-order behavior of $L_j$, optimal truncation, and exponentially small lattice sectors. The all-fixed-orders proof does not settle convergence, resurgence, or exponentially accurate counting.
\item Extend the uniform marked analysis to unequal row sums whose dimension also grows. This requires estimates beyond the exact symmetric falling-factorial cancellation used here.
\end{enumerate}

\begin{thebibliography}{9}
\bibitem{oeisdiag} OEIS Foundation, \emph{A262810: Number of alignments of $n$ sequences of length $n$}. Formula by V. Kot\v{e}\v{s}ovec, 23 March 2016. Inspected 3 October 2026. \url{https://oeis.org/A262810}.
\bibitem{oeisarray} OEIS Foundation, \emph{A262809: Number of alignments of $k$ sequences of length $n$}. Includes the positive-sum formula attributed to Peter Bala (2018). Inspected 3 October 2026. \url{https://oeis.org/A262809}.
\bibitem{ghow} J. R. Griggs, P. Hanlon, A. M. Odlyzko, and M. S. Waterman, On the number of alignments of $k$ sequences, \emph{Graphs and Combinatorics} \textbf{6} (1990), 133--146. \url{https://doi.org/10.1007/BF01787724}. Author-hosted scan: \url{https://dornsife.usc.edu/msw/wp-content/uploads/sites/236/2023/09/msw-096.pdf}.
\bibitem{pw} R. Pemantle and M. C. Wilson, Twenty combinatorial examples of asymptotics derived from multivariate generating functions, \emph{SIAM Review} \textbf{50} (2008), 199--272, Section 4.9. \url{https://www.cs.auckland.ac.nz/~mcw/Research/Outputs/PeWi2008.pdf}.
\bibitem{eger} S. Eger, On the number of many-to-many alignments of multiple sequences, \emph{Journal of Automata, Languages and Combinatorics} \textbf{20}(1) (2015), 53--65; revised arXiv version (2016), Section 5.5. \url{https://arxiv.org/abs/1511.00622}.
\bibitem{oeisnonnegative} OEIS Foundation, \emph{A316674: Nonnegative matrices with no zero columns and prescribed common row sum}. The diagonal is A316677; the recorded relation to A262809 yields the diagonal factor $2^{n-1}$. Inspected 3 October 2026. \url{https://oeis.org/A316674}. Related diagonal: \url{https://oeis.org/A316677}.
\end{thebibliography}
\end{document}
